\unnumberedchapter{第二卷参考文献（续）}
\src{205}
\noindent\textbf{合订编者说明。}第二卷此前最终版的书末源码未能恢复。以下书末资料依据本对话上传的同一法文底本重新译排；它们是原著资料的译文，不是对正文的新增数学审校。原著索引的编号与原著目录的页码均保持原义。
\begin{description}[leftmargin=3em,style=nextline]
\item[{[23]}] C. CHEVALLEY，\textit{Introduction to the theory of algebraic functions of one variable}（《单变量代数函数理论引论》），Math. Surveys，第6号，Amer. Math. Soc.，纽约，1951。
\item[{[24]}] P. JAFFARD，\textit{Les systèmes d'idéaux}（《理想系统》），巴黎，Dunod，1960。
\item[{[25]}] M. NAGATA，\textit{On the derived normal rings of Noetherian integral domains}（《论诺特整环的导出正规环》），\textit{Mem. Coll. Sci. Kyoto}，A辑，第XXIX卷（1955），293--303页。
\item[{[26]}] M. NAGATA，\textit{Existence theorems for non projective complete algebraic varieties}（《非射影完全代数簇的存在性定理》），\textit{Ill. J. Math.}，第II卷（1958），490--498页。
\end{description}

\unnumberedchapter{第二卷原著记号索引}
\src{207--208}
\begingroup\small
\begin{longtable}{@{}p{0.77\textwidth}@{\hspace{1em}}p{0.18\textwidth}@{}}
\toprule 记号与说明 & 原著条目\\\midrule\endfirsthead
\toprule 记号与说明 & 原著条目\\\midrule\endhead
\bottomrule\endfoot
$\A(X),\A(\F),\A(\B)$（$X$ 为 $S$-预概形，$\F$ 为 $\OO_X$-模，$\B$ 为 $\OO_X$-代数）&1.1.1\\
$\A(h)$（$h$ 为 $S$-态射）&1.1.2\\
$\Spec(\B)$（$\B$ 为 $\OO_S$-代数）&1.3.1\\
$\wt\M$（$\M$ 为 $\B$-模，$\B$ 为 $\OO_S$-代数）&1.4.3\\
$\mathrm T(E),\Sym(E),\Sym_A(E)$（$E$ 为 $A$-模）&1.7.1\\
$\Sym(\E),\Sym_{\A}(\E)$（$\E$ 为 $\A$-模）&1.7.4\\
$\V(\E)$（$\E$ 为拟凝聚 $\OO_S$-模）&1.7.8\\
$S_n,S_+,M_n,S^{(d)},M^{(d,k)},M^{(d)}$（$S$ 为分次环，$M$ 为分次 $S$-模）&2.1.1\\
$M(n),S(n)$（$S$ 为分次环，$M$ 为分次 $S$-模）&2.1.1\\
$\Hom_S(M,N)$（$S$ 为分次环，$M,N$ 为分次 $S$-模）&2.1.2\\
$\nrad_+,\mathfrak r_+(\frj)$（$\frj$ 为 $S_+$ 的理想）&2.1.10\\
$S_{(f)},M_{(f)}$（$f$ 为 $S_+$ 的齐次元素）&2.2.1\\
$S_{(T)},M_{(T)},S_{(\pp)},M_{(\pp)}$（$T$ 为 $S_+$ 中由齐次元素组成的乘法部分，$\pp$ 为 $S_+$ 的分次素理想）&2.2.7\\
$\Proj(S)$（$S$ 为分次环）&2.3.1\\
$V_+(E)$（$E$ 为 $S_+$ 的子集）&2.3.2\\
$D_+(f)$（$f$ 为 $S$ 的元素）&2.3.3\\
$\mathfrak j_+(Y)$（$Y$ 为 $\Proj(S)$ 的子集）&2.3.10\\
$\wt M$（$M$ 为分次 $S$-模）&2.5.2\\
$\wt u$（$u$ 为零次同态）&2.5.4\\
$\OO_X(n),\F(n)$（$X=\Proj(S)$，$\F$ 为 $\OO_X$-模）&2.5.10\\
$\lambda$（典范同态）&2.5.11\\
$\mu$（典范同态）&2.5.12\\
$X^{(d)}$（$X=\Proj(S)$，$d$ 为正整数）&2.5.16\\
$\Gamma_*(\F)$（$\F$ 为 $\OO_X$-模，$X=\Proj(S)$）&2.6.1\\
$\alpha_n,\alpha,\alpha_M$（典范同态）&2.6.2\\
$\beta,\beta_{\F}$（典范同态）&2.6.4\\
$G(\varphi),{}^a\varphi$（后一记号为滥用记号；$\varphi$ 为分次环同态）&2.8.1\\
$\Proj(\varphi)$（$\varphi$ 为分次环同态）&2.8.2\\
$\SSS^{(d)},\M^{(d,k)},\M^{(d)},\M(n)$（$\SSS$ 为分次 $\OO_Y$-代数，$\M$ 为分次 $\SSS$-模）&3.1.1\\
$\Proj(\SSS)$（$\SSS$ 为拟凝聚分次 $\OO_Y$-代数）&3.1.3\\
$X_f$（$X=\Proj(\SSS)$，$f$ 为 $\SSS_d$ 在 $Y$ 上的截面）&3.1.4\\
$\wt\M$（$\M$ 为分次 $\SSS$-模）&3.2.2\\
$\OO_X(n),\F(n)$（$X=\Proj(\SSS)$，$\F$ 为 $\OO_X$-模）&3.2.5\\
$\lambda,\mu$（典范同态）&3.2.6\\
$\Gamma_*(\F)$（$\F$ 为 $\OO_X$-模，$X=\Proj(\SSS)$）&3.3.1\\
$\alpha_n,\alpha,\alpha_{\M}$（典范同态）&3.3.2\\
$\beta,\beta_{\F}$（典范同态）&3.3.4\\
$G(\varphi),\Proj(\varphi)$（$\varphi$ 为分次 $\OO_Y$-代数同态）&3.5.1\\
$\Proj(u)$（$u$ 为从分次 $\OO_Y$-代数到分次 $\OO_{Y'}$-代数的 $g$-态射）&3.5.6\\
$r_{\LL,\psi}$（$\LL$ 为可逆 $\OO_X$-模）&3.7.1\\
$\Pb(\E),\Pb(E),\Pb_Y^n,\Pb_A^n$&4.1.1\\
$\Pb(u)$（$u$ 为 $\OO_Y$-模的满同态）&4.1.2\\
$\operatorname{Hyp}_Y(X,\E)$（$\E$ 为拟凝聚 $\OO_Y$-模）&4.2.5\\
$\varsigma$（Segre 态射）&4.3.1\\
$\F(n)$（$\F$ 为 $\OO_X$-模，$X$ 配备一个可逆 $\OO_X$-模）&4.5.1\\
$\varepsilon$（恒等自同构）&4.5.1\\
$P(u,T),\sigma_i(u)$（$u$ 为有限基自由模的自同态）&6.4.1\\
$P(u,T),\sigma_i(u),\det u$（$u$ 为整环或约化诺特环上有限型模的自同态）&6.4.2；6.4.7\\
$\sigma_i(u),\det u$（$u$ 为局部自由 $\A$-模的自同态）&6.4.8\\
$\sigma_i,\det$（层的同态）&6.4.8--6.4.10\\
$\sigma_i(u),\det u$（$u$ 为局部整预概形或局部诺特约化预概形上有限型 $\OO_X$-模的自同态）&6.4.9；6.4.10\\
$N_{\B/\A}(f)$（$f$ 为 $\A$-代数 $\B$ 的截面）&6.5.1\\
$N_{\B/\A}(\LL')$（$\LL'$ 为可逆 $\B$-模）&6.5.2\\
$N_{\B/\A}(h')$（$h'$ 为可逆 $\B$-模的同态）&6.5.3\\
$N_{X'/X}(\LL')$（$X'$ 在 $X$ 上有限，$\LL'$ 为可逆 $\OO_{X'}$-模）&6.5.5\\
$N_{X'/X}(h')$（$h'$ 为可逆 $\OO_{X'}$-模的同态）&6.5.5\\
$\SSS_X$（$X=\Proj(\SSS)$，$\SSS$ 为 $\OO_Y$ 的理想层之直和）&8.1.5\\
$S^{\geq},S^{\leq},S_f^{\geq},S_f^{\leq},M^{\geq},M^{\leq},M_f^{\geq},M_f^{\leq}$（$S$ 为分次环，$M$ 为分次 $S$-模，$f$ 为 $S_+$ 的齐次元素）&8.2.1\\
$\widehat S$（$S$ 为分次环）&8.2.2\\
$\widehat M$（$M$ 为分次 $S$-模）&8.2.4\\
$S_{[n]},S^{\natural},f^{\natural}$（$S$ 为分次环，$f$ 为其中的齐次元素）&8.2.6\\
$M_{[n]},M^{\natural}$（$M$ 为分次 $S$-模）&8.2.8\\
$\widehat\SSS$（$\SSS$ 为分次 $\OO_Y$-代数）&3.8.1\textsuperscript{*}\\
$\mathbf G_m$（群概形）&8.3.9\\
$\SSS_X$（$X=\Proj(\SSS)$，$\SSS$ 为分次 $\OO_Y$-代数）&8.6.1\\
$\SSS_X^{\geq},C_X,\widehat C_X,E_X,\widehat E_X$（$X=\Proj(\SSS)$，$\SSS_0=\OO_Y$）&8.6.1\\
$\SSS_{[n]},\SSS^{\natural}$（$\SSS$ 为分次 $\OO_Y$-代数）&8.7.2\\
$\sProj_0(\M),\sProj(\M),\M_X,\widehat\M,\M^\square$（$\M$ 为分次 $\SSS$-模）&8.12.1\\
$\M_X^{\geq}$（$X=\Proj(\SSS)$，$\M$ 为分次 $\SSS$-模）&8.12.5\\
$\M_{[n]},\M^{\natural}$（$\M$ 为分次 $\SSS$-模）&8.12.9\\
$(\wt\N)^-$（$\N$ 为分次 $\SSS$-模的 $\SSS$-子模）&8.13.1\\
$\overline\N$（$\N$ 为分次 $\SSS$-模的 $\SSS$-子模）&8.13.2\\
$\boldsymbol\Gamma_*(\F)$（$X=\Proj(\SSS)$，$\F$ 为分次 $\SSS_X$-模）&8.14.5\\
$\SSS_{X'},\boldsymbol\Gamma'_*(\F'),\sProj'(\M),\alpha',\beta'$（$X'$ 为 $X=\Proj(\SSS)$ 的开集，$\F'$ 为分次 $\SSS_{X'}$-模，$\M$ 为分次 $\SSS$-模）&8.14.6\\
\end{longtable}
\endgroup

\noindent\textbf{译注。}\textsuperscript{*}原著记号索引此处印为3.8.1；相应记号见正文8.3.1。索引原编号予以保留。

\unnumberedchapter{第二卷原著术语索引}
\src{209--211}
\begingroup\small
\begin{longtable}{@{}p{0.77\textwidth}@{\hspace{1em}}p{0.18\textwidth}@{}}
\toprule 术语 & 原著条目\\\midrule\endfirsthead
\toprule 术语 & 原著条目\\\midrule\endhead\bottomrule\endfoot
仿射态射 &1.6.1\\
一个预概形之上的仿射预概形 &1.2.1\\
分次环上的分次代数 &2.1.2\\
$A$-模的对称代数 &1.7.1\\
$\A$-模的对称 $\A$-代数 &1.7.4\\
本质约化的分次 $\OO_Y$-代数 &3.1.12\\
整分次 $\OO_Y$-代数 &3.1.12\\
丰沛可逆 $\OO_X$-模 &4.5.3\\
相对于 $f$ 丰沛、$f$-丰沛、$Y$-丰沛的可逆 $\OO_X$-模 &4.6.1\\
本质整分次环 &2.1.11\\
本质约化分次环 &2.1.10\\
与 $\B$-模相伴的层 &1.4.3\\
与分次 $S$-模相伴的层 &2.5.3\\
与分次 $\SSS$-模相伴的层 &3.2.2\\
与 $\OO_S$-代数相伴的 $S$-概形 &1.3.1\\
（TF）条件、（TN）条件 &2.7.2\\
（TF）条件、（TN）条件 &3.4.2\\
由分次 $\OO_Y$-代数定义的仿射锥、射影锥 &8.3.1\\
由分次 $\OO_Y$-代数定义的去顶点仿射锥、去顶点射影锥 &8.3.4\\
预概形 $\Proj(\SSS)$ 的仿射投影锥、射影投影锥 &8.3.1\\
域上的代数曲线 &7.4.2\\
完全代数曲线 &7.4.7\\
Serre 判别准则 &5.2.1\\
整环上有限型模自同态的行列式 &6.4.2\\
约化诺特环上有限型模自同态的行列式 &6.4.7\\
局部自由 $\A$-模自同态的行列式 &6.4.8\\
局部整预概形 $X$ 上有限型 $\OO_X$-模自同态的行列式 &6.4.9\\
局部诺特约化预概形 $X$ 上有限型 $\OO_X$-模自同态的行列式 &6.4.10\\
爆破预概形 &8.1.3\\
整态射 &6.1.1\\
一个预概形之上的整预概形 &6.1.1\\
整拟凝聚 $\OO_S$-代数 &6.1.2\\
$\A$-代数中在 $\A$ 上整的截面 &6.3.1\\
射影丛的基本层 &4.1.1\\
环层 $\A$ 在一个 $\A$-代数中的整闭包 &6.3.2\\
预概形 $X$ 关于一个 $\OO_X$-代数的整闭包 &6.3.4\\
仿射锥的射影闭包 &8.3.1\\
射影丛在 $k(y)$ 的扩张 $K$ 上有理的几何纤维 &4.2.6\\
由 $\OO_Y$-模定义的射影丛 &4.1.1\\
由 $\OO_S$-模定义的向量丛 &1.7.8\\
有限态射 &6.1.1\\
一个预概形之上的有限预概形 &6.1.1\\
有限拟凝聚 $\OO_S$-代数 &6.1.2\\
局部自由模自同态的初等对称函数 &6.4.1\\
整环上有限型模自同态的初等对称函数 &6.4.2\\
约化诺特环上有限型模自同态的初等对称函数 &6.4.7\\
常秩局部自由 $\A$-模自同态的初等对称函数 &6.4.8\\
局部整预概形上有限型 $\OO_X$-模自同态的初等对称函数 &6.4.9\\
局部诺特约化预概形上有限型 $\OO_X$-模自同态的初等对称函数 &6.4.10\\
一个 $\A$-代数的截面的初等对称函数 &6.5.1\\
分次 $\SSS$-模的 $\SSS$-子模的齐次化 &8.13.2\\
分次环同态 &2.1.2\\
分次模之间的 $k$ 次同态 &2.1.2\\
$S_+$ 的理想、$S_+$ 的分次素理想（$S$ 为分次环）&2.1.10\\
Chow 引理 &5.6.1\\
射影锥的无穷远处 &8.3.3\\
典范态射 $G(\LL)\to\Proj(S)$，其中 $S=\bigoplus_{n\geq0}\Gamma(X,\LL^{\otimes n})$ &4.5.1\\
典范态射 $X\to\Spec(\Gamma(X,\OO_X))$ &5.1.1\\
Segre 态射 &4.3.1\\
$S_+$ 的幂零根（$S$ 为分次环）&2.1.10\\
$\SSS_+$ 的幂零根（$\SSS$ 为拟凝聚分次 $\OO_Y$-代数）&3.1.12\\
一个 $\A$-代数的截面的范数 &6.5.1\\
可逆 $\B$-模的范数 &6.5.2.3\\
可逆 $\B$-模的截面的范数 &6.5.3\\
可逆 $\OO_{X'}$-模的范数 &6.5.5\\
周期分次代数 &8.14.12\\
整环上有限型模自同态的特征多项式 &6.4.2\\
约化诺特环上有限型模自同态的特征多项式 &6.4.7\\
仿射锥（射影锥）的顶点预概形 &8.3.3\\
有限表示的分次模 &2.1.1\\
有限表示的分次 $\SSS$-模 &3.1.1\\
两个分次模的分次张量积 &2.1.2\\
射影态射 &5.5.2\\
一个预概形之上的射影预概形 &5.5.2\\
固有态射 &5.4.1\\
固有子集 &5.4.10\\
一个预概形之上的固有预概形 &5.4.1\\
拟仿射态射 &5.1.1\\
一个预概形之上的拟仿射预概形 &5.1.1\\
拟仿射概形 &5.1.1\\
拟有限态射 &6.2.3\\
一个预概形之上的拟有限预概形 &6.2.3\\
拟射影态射 &5.3.1\\
一个预概形之上的拟射影预概形 &5.3.1\\
$S_+$ 中理想的根（$S$ 为分次环）&2.1.10\\
去顶点射影锥到其无穷远处的典范收缩 &8.3.5\\
$\OO_X(-1)$ 的典范截面（$X$ 为爆破预概形）&8.1.9\\
仿射锥的零截面 &8.3.3\\
向量丛的零截面 &1.7.9\\
仿射锥（射影锥）的顶点截面 &8.3.3\\
一个 $\OO_S$-代数的谱 &1.3.1\\
分次环的齐次素谱 &2.3.1\\
拟凝聚分次 $\OO_S$-代数的齐次谱 &3.1.3\\
分次模同态的（TN）-单射、（TN）-满射、（TN）-双射 &2.7.2\\
分次模的（TN）-同构 &2.7.2\\
分次环同态的（TN）-单射、（TN）-满射、（TN）-双射 &2.9.1\\
分次环的（TN）-同构 &2.9.1\\
分次 $\SSS$-模同态的（TN）-单射、（TN）-满射、（TN）-双射 &3.4.2\\
分次 $\SSS$-模的（TN）-同构 &3.4.2\\
分次 $\OO_Y$-代数同态的（TN）-单射、（TN）-满射、（TN）-双射 &3.6.1\\
分次 $\OO_Y$-代数的（TN）-同构 &3.6.1\\
$\Proj(S)$ 上的谱拓扑 &2.3.3\\
相对于 $q$ 甚丰沛、相对于 $Y$ 甚丰沛的可逆 $\OO_X$-模 &4.4.2\\
有限型分次 $\SSS$-模 &3.1.1\\
普遍闭态射 &5.4.9\\
\end{longtable}
\endgroup

\unnumberedchapter{第二卷原著目录}
\src{213--215}
\noindent 下表保留法文原著的印刷页码，不是本合订本的页码。
\begingroup\small
\begin{longtable}{@{}p{0.85\textwidth}@{\hspace{1em}}r@{}}
\toprule 内容 & 原书页码\\\midrule\endfirsthead
\toprule 内容 & 原书页码\\\midrule\endhead\bottomrule\endfoot
\textbf{第II章　若干类态射的初等整体研究} &5\\
\textbf{§1　仿射态射}&5\\
1.1　$S$-预概形与 $\OO_S$-代数&5\\
1.2　一个预概形之上的仿射预概形&6\\
1.3　与一个 $\OO_S$-代数相伴的、在 $S$ 之上仿射的预概形&8\\
1.4　在 $S$ 之上仿射的预概形上的拟凝聚层&9\\
1.5　底预概形的变换&12\\
1.6　仿射态射&14\\
1.7　与一个模层相伴的向量丛&14\\
\textbf{§2　齐次素谱}&19\\
2.1　分次环与分次模的一般性质&19\\
2.2　分次环的分式环&23\\
2.3　分次环的齐次素谱&25\\
2.4　$\Proj(S)$ 上的概形结构&28\\
2.5　与一个分次模相伴的层&30\\
2.6　与 $\Proj(S)$ 上的层相伴的分次 $S$-模&36\\
2.7　有限性条件&38\\
2.8　函子性质&41\\
2.9　概形 $\Proj(S)$ 的闭子预概形&48\\
\textbf{§3　分次代数层的齐次谱}&49\\
3.1　拟凝聚分次 $\OO_Y$-代数的齐次谱&49\\
3.2　与分次 $\SSS$-模相伴的 $\Proj(\SSS)$ 上的层&54\\
3.3　与 $\Proj(\SSS)$ 上的层相伴的分次 $\SSS$-模&56\\
3.4　有限性条件&59\\
3.5　函子性质&61\\
3.6　预概形 $\Proj(\SSS)$ 的闭子预概形&64\\
3.7　从预概形到齐次谱的态射&65\\
3.8　浸入齐次谱的判别准则&69\\
\textbf{§4　射影丛。丰沛层}&71\\
4.1　射影丛的定义&71\\
4.2　从预概形到射影丛的态射&72\\
4.3　Segre 态射&76\\
4.4　到射影丛的浸入。甚丰沛层&78\\
4.5　丰沛层&83\\
4.6　相对丰沛层&89\\
\textbf{§5　拟仿射态射；拟射影态射；固有态射；射影态射}&94\\
5.1　拟仿射态射&94\\
5.2　Serre 判别准则&97\\
5.3　拟射影态射&99\\
5.4　固有态射与普遍闭态射&100\\
5.5　射影态射&103\\
5.6　Chow 引理&106\\
\textbf{§6　整态射与有限态射}&110\\
6.1　在另一预概形上为整的预概形&110\\
6.2　拟有限态射&114\\
6.3　预概形的整闭包&116\\
6.4　模层自同态的行列式&120\\
6.5　可逆层的范数&125\\
6.6　应用：丰沛性的判别准则&130\\
6.7　Chevalley 定理&135\\
\textbf{§7　赋值判别准则}&138\\
7.1　赋值环的回顾&138\\
7.2　分离性的赋值判别准则&141\\
7.3　固有性的赋值判别准则&143\\
7.4　代数曲线与一维函数域&148\\
\textbf{§8　爆破概形；投影锥；射影闭包}&152\\
8.1　爆破预概形&152\\
8.2　分次环局部化的初步结果&157\\
8.3　投影锥&162\\
8.4　向量丛的射影闭包&168\\
8.5　函子性质&169\\
8.6　去顶点锥的一个典范同构&171\\
8.7　投影锥的爆破&173\\
8.8　丰沛层与收缩&177\\
8.9　Grauert 丰沛性判别准则：陈述&182\\
8.10　Grauert 丰沛性判别准则：证明&184\\
8.11　收缩的唯一性&189\\
8.12　投影锥上的拟凝聚层&191\\
8.13　子层与闭子概形的射影闭包&195\\
8.14　与分次 $\SSS$-模相伴的层的补充&197\\
参考文献（续）&205\\
记号索引&207\\
术语索引&209\\
勘误与增补（清单一）&217\\
\end{longtable}\endgroup
\begin{flushright}1960年7月22日收到。\end{flushright}

\unnumberedchapter{勘误与增补（清单一）}
\src{217--222}
\noindent\textbf{合订编者说明。}本清单是原作者随第二卷发表的、针对第0章和第I章的勘误与增补；其中的页码均指第一卷原版。保留其独立位置，不将这些改动无说明地并入第一卷正文。以下条目号及公式号按清单原文保留。

\section*{第0章}
\addcontentsline{toc}{section}{勘误与增补：第0章}
\begin{description}[style=nextline,leftmargin=2em]
\item[（0，2.1.3）] 第21页倒数第15行，将 \textit{inter-section} 改为 \textit{intersection}。
\item[（0，2.1.6）] 第22页第13行，将 $Z_i\cap\complement(\bigcup_{j\ne i}Z_j)$ 改为 $\bigcup_{j\ne i}Z_j$。\transnote{此处按勘误清单的可见字面转录。原段落中涉及开集与补集的定义，不能据此字面条目无说明地自动替换正文。}
\item[（0，3.1.6）] 第25页倒数第17行，将 $\Gamma(U,\F)$ 改为 $\Gamma(V,\F)$。
\item[（0，4.1.1）] 第36页第12行，将 $\G\to\F$ 改为 $\B\to\A$。
\item[（0，4.2.1）] 第39页倒数第11行，将 \textit{faisceaux}（层，复数）改为 \textit{faisceau}（层，单数）。
\item[（0，5.1.2）] 在第45页第14行之后加入：若 $\G$ 是由其在 $Y$ 上的截面生成的 $\OO_Y$-模，则 $f^*(\G)$ 由其在 $X$ 上的截面生成，因为 $f^*$ 是右正合函子。
\item[（0，5.2.4）] 第46页倒数第18行，将 $f_U$ 改为 $f_U^*$。
\item[（0，5.3.4）] 在第47页倒数第16行之前加入：若 $\A\to\B\to\C\to\mathscr D\to\E$ 是 $\OO_X$-模正合列，且 $\A,\B,\mathscr D,\E$ 都凝聚，则 $\C$ 凝聚。
\item[（0，5.4.1）] 在第49页第2行之后加入：假设 $\OO_X$ 凝聚，且 $\F$ 为凝聚 $\OO_X$-模。若在点 $x\in X$，$\F_x$ 是秩为 $n$ 的自由 $\OO_x$-模，则存在 $x$ 的邻域 $U$，使 $\F|U$ 局部自由且秩为 $n$。事实上，此时 $\F_x\simeq\OO_x^n$，由（5.2.7）即得此命题。
\item[（0，6.6.2）] 将第59页第10、11行替换为：“事实上，这由（6.4.1，d））得到。”
\item[（0，6.7.1）] 在第59页倒数第10行之前加入：若 $f$ 为平坦态射，也称 $X$ 在 $Y$ 上平坦，或称 $X$ 是 $Y$-平坦的。
\item[（0，7.2.9）] 第65页倒数第11行，将 $M_n$ 改为 $M_{n-1}$。
\end{description}

\section*{第I章}
\addcontentsline{toc}{section}{勘误与增补：第I章}
\begin{description}[style=nextline,leftmargin=2em]
\item[（I，1.2.1）] 第83页第18行，两处 $K$ 均改为 $k$。
\item[（I，1.2.7）] 第84页倒数第16行，将 $A$（两处）改为 $A'$，将 $\nrad$ 改为 $\nrad'$；倒数第17行，将 $X$ 改为 $X'$。
\item[（I，1.7.3）及（I，2.2.4）] 根据 J. Tate 的一个观察，这些条目中的陈述可推广如下。
\end{description}
\subsection*{1.8　从局部赋环空间到仿射概形的态射}
\addcontentsline{toc}{subsection}{增补 I，1.8　从局部赋环空间到仿射概形的态射}
\erritem{1.8.1}{命题}设 $(S,\OO_S)$ 为仿射概形，$(X,\OO_X)$ 为局部赋环空间。则环同态
$\Gamma(S,\OO_S)\to\Gamma(X,\OO_X)$ 的集合，与满足下述条件的赋环空间态射
$(\psi,\theta):(X,\OO_X)\to(S,\OO_S)$ 的集合之间有典范双射：对每个 $x\in X$，
$\theta_x^\sharp:\OO_{\psi(x)}\to\OO_x$ 是局部同态。
\begin{proof}
首先，若 $(X,\OO_X)$、$(S,\OO_S)$ 为任意两个赋环空间，则从前者到后者的赋环空间态射 $(\psi,\theta)$ 典范地确定一个环同态
$\Gamma(\theta):\Gamma(S,\OO_S)\to\Gamma(X,\OO_X)$，因而得到第一个映射
\begin{equation}
\rho:\Hom((X,\OO_X),(S,\OO_S))\longrightarrow
\Hom(\Gamma(S,\OO_S),\Gamma(X,\OO_X)).\tag{1.8.1.1}
\end{equation}
反之，在命题的假设下，令 $A=\Gamma(S,\OO_S)$，考虑环同态 $\varphi:A\to\Gamma(X,\OO_X)$。对每个 $x\in X$，满足 $\varphi(f)(x)=0$ 的 $f\in A$ 所成集合显然是 $A$ 的素理想，因为 $\OO_x/\mm_x=k(x)$ 是域。因此它是 $S=\Spec(A)$ 的一个元素，仍记为 ${}^a\varphi(x)$。此外，对每个 $f\in A$，依定义（0，5.5.2），
$({}^a\varphi)^{-1}(D(f))=X_f$，这证明 ${}^a\varphi:X\to S$ 连续。

再定义 $\OO_S$-模同态 $\wt\varphi:\OO_S\to{}^a\varphi_*(\OO_X)$。对每个 $f\in A$，有 $\Gamma(D(f),\OO_S)=A_f$（1.3.6）；对每个 $s\in A$，将 $s/f\in A_f$ 对应于
\[
(\varphi(s)|X_f)(\varphi(f)|X_f)^{-1}\in
\Gamma(X_f,\OO_X)=\Gamma(D(f),{}^a\varphi_*(\OO_X)).
\]
通过从 $D(f)$ 换为 $D(fg)$，立即可验证这确实定义一个 $\OO_S$-模同态，从而得到赋环空间态射 $({}^a\varphi,\wt\varphi)$。仍沿用这些记号，为简便令 $y={}^a\varphi(x)$，立即有（0，3.7.1）
\[
\wt\varphi_x^\sharp(s_y/f_y)=(\varphi(s)_x)(\varphi(f)_x)^{-1}.
\]
按定义，$s_y\in\mm_y$ 等价于 $\varphi(s)_x\in\mm_x$，故 $\wt\varphi_x^\sharp:\OO_y\to\OO_x$ 是局部同态。由此定义第二个映射
$\sigma:\Hom(\Gamma(S,\OO_S),\Gamma(X,\OO_X))\to\Omega$，其中 $\Omega$ 为满足每个 $\theta_x^\sharp$ 都局部的态射 $(\psi,\theta):(X,\OO_X)\to(S,\OO_S)$ 所成的集合。

余下须证明 $\sigma$ 与 $\rho$ 在 $\Omega$ 上的限制互逆。由 $\wt\varphi$ 的定义立即有 $\Gamma(\wt\varphi)=\varphi$，故 $\rho\circ\sigma$ 为恒等映射。为证明 $\sigma\circ\rho$ 为恒等映射，取 $(\psi,\theta)\in\Omega$ 并令 $\varphi=\Gamma(\theta)$。由 $\theta_x^\sharp$ 的假设，通过取商得到域的单同态
$\theta_x^\natural:k(\psi(x))\to k(x)$，满足对每个 $f\in A=\Gamma(S,\OO_S)$，
$\theta_x^\natural(f(\psi(x)))=\varphi(f)(x)$。因此 $f(\psi(x))=0$ 等价于 $\varphi(f)(x)=0$，已证明 ${}^a\varphi=\psi$。另一方面，由定义，图
\[
\begin{tikzcd}
A\arrow[r,"\varphi"]\arrow[d]&\Gamma(X,\OO_X)\arrow[d]\\
A_{\psi(x)}\arrow[r,"\theta_x^\sharp"']&\OO_x
\end{tikzcd}
\]
交换；以 $\wt\varphi_x^\sharp$ 代替 $\theta_x^\sharp$ 所得的类似图也交换，故（0，1.2.4）$\wt\varphi_x^\sharp=\theta_x^\sharp$，从而 $\wt\varphi=\theta$。
\end{proof}

\erritem{1.8.2}{}当 $(X,\OO_X)$、$(Y,\OO_Y)$ 为局部赋环空间时，我们将考虑满足对每个 $x\in X$，$\theta_x^\sharp:\OO_{\psi(x)}\to\OO_x$ 为局部同态的态射 $(\psi,\theta):(X,\OO_X)\to(Y,\OO_Y)$。今后说到局部赋环空间的态射时，均指这种态射。用这一态射定义，局部赋环空间显然组成一个范畴；对此范畴的两个对象 $X,Y$，$\Hom(X,Y)$ 表示从 $X$ 到 $Y$ 的局部赋环空间态射集合，即（1.8.1）中记作 $\Omega$ 的集合。需要考虑从 $X$ 到 $Y$ 的全部赋环空间态射时，为免混淆，记为 $\Hom_{\mathrm{an}}(X,Y)$。于是（1.8.1.1）写作
\begin{equation}
\rho:\Hom_{\mathrm{an}}(X,Y)\longrightarrow
\Hom(\Gamma(Y,\OO_Y),\Gamma(X,\OO_X)),\tag{1.8.2.1}
\end{equation}
其限制
\begin{equation}
\rho':\Hom(X,Y)\longrightarrow
\Hom(\Gamma(Y,\OO_Y),\Gamma(X,\OO_X))\tag{1.8.2.2}
\end{equation}
是局部赋环空间范畴中关于 $X,Y$ 的函子性映射。

\erritem{1.8.3}{推论}设 $(Y,\OO_Y)$ 为局部赋环空间。$Y$ 是仿射概形，当且仅当对于每个局部赋环空间 $(X,\OO_X)$，映射（1.8.2.2）为双射。
\begin{proof}
命题（1.8.1）说明必要性。反之，假设此条件成立，令 $A=\Gamma(Y,\OO_Y)$。由假设和（1.8.1），从局部赋环空间范畴到集合范畴的函子 $X\mapsto\Hom(X,Y)$ 和 $X\mapsto\Hom(X,\Spec(A))$ 同构。已知这蕴含一个典范同构 $Y\simeq\Spec(A)$ 的存在（参见0，8）。\transnote{原版末句写作 $X\simeq\Spec(A)$；其所指的固定对象是 $Y$，故在此说明该字母误植。}
\end{proof}

\erritem{1.8.4}{}设 $S=\Spec(A)$ 为仿射概形。以 $(S',A')$ 表示底空间只有一个点、结构层 $A'$ 为环 $A$ 在 $S'$ 上定义的层（必为常值层）的赋环空间。令 $\pi:S\to S'$ 为唯一的映射。另一方面，对 $S$ 的每个开集 $U$，有典范映射
$\Gamma(S',A')=\Gamma(S,\OO_S)\to\Gamma(U,\OO_S)$；它定义环层的 $\pi$-态射 $\iota:A'\to\OO_S$。这样典范地定义赋环空间态射 $i=(\pi,\iota):(S,\OO_S)\to(S',A')$。对每个 $A$-模 $M$，以 $M'$ 表示由 $M$ 在 $S'$ 上定义的常值层，它显然是一个 $A'$-模。显然 $i_*(\wt M)=M'$（1.3.7）。

\erritem{1.8.5}{引理}采用（1.8.4）的记号，对每个 $A$-模 $M$，典范函子性 $\OO_S$-同态（0，4.4.3.3）
\begin{equation}i^*(i_*(\wt M))\longrightarrow\wt M\tag{1.8.5.1}\end{equation}
为同构。
\begin{proof}
（1.8.5.1）的两端都右正合（函子 $M\mapsto i_*(\wt M)$ 显然正合），并与直和交换。将 $M$ 看成同态 $A^{(I)}\to A^{(J)}$ 的余核，可归结到 $M=A$ 的情形；在此情形结论显然。
\end{proof}

\erritem{1.8.6}{推论}设 $(X,\OO_X)$ 为赋环空间，$u:X\to S$ 为赋环空间态射。对每个 $A$-模 $M$，采用（1.8.4）的记号，有 $\OO_X$-模的典范函子性同构
\begin{equation}u^*(\wt M)\xrightarrow{\sim}u^*(i^*(M')).\tag{1.8.6.1}\end{equation}

\erritem{1.8.7}{推论}在（1.8.6）的假设下，对每个 $A$-模 $M$ 和每个 $\OO_X$-模 $\F$，有关于 $M,\F$ 的典范函子性同构
\begin{equation}
\Hom_{\OO_S}(\wt M,u_*(\F))\xrightarrow{\sim}
\Hom_A(M,\Gamma(X,\F)).\tag{1.8.7.1}
\end{equation}
事实上，由（0，4.4.3）及引理（1.8.5），有双函子的典范同构
\[
\Hom_{\OO_S}(\wt M,u_*(\F))\simeq
\Hom_{A'}(M',i_*(u_*(\F))).
\]
右端显然就是 $\Hom_A(M,\Gamma(X,\F))$。注意，（1.8.7.1）将每个 $\OO_S$-同态 $h:\wt M\to u_*(\F)$，也就是每个 $u$-态射 $\wt M\to\F$，映到 $A$-同态
$\Gamma(h):M\to\Gamma(S,u_*(\F))=\Gamma(X,\F)$。

\erritem{1.8.8}{}采用（1.8.4）的记号，由（0，4.1.1），每个赋环空间态射 $(\psi,\theta):X\to S'$ 显然等价于给定环同态 $A\to\Gamma(X,\OO_X)$。所以也可将（1.8.1）解释为定义一个典范双射
$\Hom(X,S)\simeq\Hom(X,S')$，其中右端当然指赋环空间态射，因为一般 $A$ 并非局部环。更一般地，若 $X,Y$ 为局部赋环空间，$(Y',A')$ 为底空间只有一点且环层 $A'$ 是由 $\Gamma(Y,\OO_Y)$ 定义的常值层的赋环空间，则可将（1.8.2.1）解释为映射
\begin{equation}\rho:\Hom_{\mathrm{an}}(X,Y)\to\Hom(X,Y').\tag{1.8.8.1}\end{equation}
因而（1.8.3）的结果说明，局部赋环空间中的仿射概形可刻画为满足如下条件的对象：对每个局部赋环空间 $X$，$\rho$ 在 $\Hom(X,Y)$ 上的限制
\begin{equation}\rho':\Hom(X,Y)\to\Hom(X,Y')\tag{1.8.8.2}\end{equation}
为双射。在以后的一个章节中，我们将推广这一定义，使之能够将一个局部赋环空间与任意赋环空间 $Z$ 相联系，而不再只限于底空间只有一点的赋环空间；我们仍将它称为 $\Spec(Z)$。这将成为任意赋环空间之上预概形的“相对”理论的出发点，并推广第I章的结果。

\erritem{1.8.9}{}可以把局部赋环空间 $X$ 与 $\OO_X$-模 $\F$ 构成的二元组 $(X,\F)$ 看作一个范畴的对象。其中态射 $(X,\F)\to(Y,\G)$ 是一个二元组 $(u,h)$，由局部赋环空间态射 $u:X\to Y$ 和模的 $u$-态射 $h:\G\to\F$ 构成。固定 $(X,\F),(Y,\G)$ 时，这些态射所成集合记为 $\Hom((X,\F),(Y,\G))$。映射 $(u,h)\mapsto(\rho'(u),\Gamma(h))$ 是典范映射
\begin{equation}
\begin{aligned}
\Hom((X,\F),(Y,\G))\longrightarrow
\Hom\bigl(& (\Gamma(Y,\OO_Y),\Gamma(Y,\G)),\\
& (\Gamma(X,\OO_X),\Gamma(X,\F))\bigr),
\end{aligned}\tag{1.8.9.1}
\end{equation}
它关于 $(X,\F),(Y,\G)$ 是函子性的；右端为与所考虑的环和模相应的双同态集合（0，1.0.2）。

\erritem{1.8.10}{推论}设 $Y$ 为局部赋环空间，$\G$ 为 $\OO_Y$-模。$Y$ 是仿射概形且 $\G$ 是拟凝聚 $\OO_Y$-模，当且仅当对于每个由局部赋环空间 $X$ 和 $\OO_X$-模 $\F$ 构成的二元组 $(X,\F)$，典范映射（1.8.8.1）为双射。\transnote{此处保留勘误清单中的引文编号；按所述二元组及映射的定义，其应指（1.8.9.1）。}

将这一论证留给读者；它仿照（1.8.3）的论证，并使用（1.8.1）与（1.8.7）。

\erritem{1.8.11}{注}（1.7.3）、（1.7.4）、（2.2.4）的陈述及定义（1.6.1）都是（1.8.1）的特殊情形；同样，（2.2.5）由（1.8.7）得到。命题（1.8.7）也推出（1.6.3），因而还推出（1.6.4），作为其特殊情形。事实上，若 $X$ 仿射且 $\Gamma(X,\F)=N$，则函子
\[
M\longmapsto\Hom_{\OO_S}(\wt M,u_*(\wt N)),\qquad
M\longmapsto\Hom_{\OO_S}(\wt M,(N_{[\varphi]})^{\sim})
\]
由（1.8.7）及（1.3.8）同构，其中 $\varphi:A\to\Gamma(X,\OO_X)$ 对应于 $u$。最后，（1.6.5），从而（1.6.6），由（1.8.6）以及以下事实得到：对每个 $f\in A$，两个 $A_f$-模
$N'\otimes_{A'}A_f$ 和 $(N'\otimes_{A'}A)_f$ 典范同构，记号如（1.6.5）。

\medskip\noindent 在（I，3.2.8）之后插入：
\erritem{3.2.9}{}由（1.8.1），可将（3.2.2）精确化如下：$Z=\Spec(B\otimes_A C)$ 不仅是 $X=\Spec(B)$ 与 $Y=\Spec(C)$ 在 $S$-预概形范畴中的积，也是它们在 $S$ 上的局部赋环空间范畴中的积，其中 $S$-态射仿照（2.5.2）定义。（3.2.6）的证明因而实际上表明，对任意两个 $S$-预概形 $X,Y$，预概形 $X\times_S Y$ 不仅是它们在 $S$-预概形范畴中的积，也是它们在预概形 $S$ 上的局部赋环空间范畴中的积。

\begin{description}[style=nextline,leftmargin=2em]
\item[（I，4.4.5）] 第126页倒数第14行，将 $B$ 改为 $A$；倒数第13行，将“$A$-代数”改为“$B$-代数”。
\item[（I，5.3.5）] 第132页倒数第2行，将 $Y\to S$ 改为 $g:Y\to S$。
\item[（I，6.3.2.1）] 第144页倒数第6行，应读作 $D(g_i)\subset W$。
\item[（I，7.3.8）] 以如下文字替换原有的错误文字：
\end{description}
设 $X,Y$ 为两个整预概形，因此 $\R(X)$（分别地，$\R(Y)$）是拟凝聚 $\OO_X$-模（分别地，$\OO_Y$-模）（7.3.3）。设 $f:X\to Y$ 为占优态射；则有 $\OO_X$-模的典范同态
\begin{equation}\tau:f^*(\R(Y))\longrightarrow\R(X).\tag{7.3.8.1}\end{equation}
先设 $X=\Spec(A)$、$Y=\Spec(B)$ 仿射，$A,B$ 为整环，$f$ 对应于单同态 $B\to A$；它延拓为从 $B$ 的分式域 $L$ 到 $A$ 的分式域 $K$ 的单同态 $L\to K$。于是（7.3.8.1）对应于典范同态 $L\otimes_B A\to K$（1.6.5）。一般地，对每对满足 $f(U)\subset V$ 的非空仿射开集 $U\subset X,V\subset Y$，按上述方式定义同态 $\tau_{U,V}$。立即可见，若 $U'\subset U,V'\subset V,f(U')\subset V'$，则 $\tau_{U,V}$ 延拓 $\tau_{U',V'}$，由此得到所需断言。

若 $x,y$ 分别为 $X,Y$ 的泛点，则 $f(x)=y$，并且
\[
(f^*(\R(Y)))_x=\OO_y\otimes_{\OO_y}\OO_x=\OO_x
\]
（0，4.3.1），因此 $\tau_x$ 是同构。
\begin{description}[style=nextline,leftmargin=2em]
\item[（I，9.5.2）] 在第176页最后3行和第177页最前5行中，将所有 $X'$ 改为 $Y'$，所有 $X''$ 改为 $Y''$。
\item[（I，10.5.4）] 第187页倒数第11行，将（10.3.4）改为（10.3.5）。
\item[（I，10.6.3）] 第189页倒数第3行，将 $X$ 改为 $\frX$。
\item[（I，10.14.2）] 第210页第5行，将“凝聚 $\OO_{\frX}$-模”改为“$\OO_{\frX}$ 的凝聚理想”；第7行，将“凝聚 $\OO_{\frX}$-模”改为“$\OO_{\frX}$ 的凝聚理想”。
\end{description}
